% Parte: teoria de conjuntos
% Capítulo: ordinais
% Ficheiro: introduction

\documentclass[../../../include/open-logic-section]{subfiles}

\begin{document}

\olfileid[pt]{sth}{ordinals}{intro}

\olsection{Introdução}

No \olref[z][]{chap}, postulamos que existe um estágio infinito
da hierarquia, na forma de \stagesinf{} (ver também o nosso axioma do
Infinito). Contudo, dado \stagessucc{}, não podemos parar no estágio
infinito; temos de continuar. Assim: no estágio seguinte ao primeiro estágio
infinito, formamos todas as coleções possíveis de conjuntos
que estavam disponíveis no primeiro estágio infinito; e repetimos; e
repetimos; e repetimos; \dots

Implicitamente, o que aqui aconteceu foi começarmos a invocar uma noção
``intuitiva'' de número, segundo a qual podem existir números
\emph{depois} de todos os números naturais. Em particular, a noção em
causa é a de \emph{ordinal transfinito}. O objetivo deste capítulo é
tornar esta ideia mais rigorosa. Exploraremos a noção geral de ordinal e
depois definiremos explicitamente certos conjuntos para serem os nossos
ordinais. 

\end{document}
